(1) 设特征坐标 \(\xi=x-ct,\eta=x+ct\)，波动算子分解为 \(\partial_t^2-c^2\partial_x^2=-4c^2\partial_\xi\partial_\eta\)。因此通解是 \(y(x,t)=F(x-ct)+G(x+ct)\)。由初始条件联立得到 \(y(x,t)=[f(x-ct)+f(x+ct)]/2+(1/(2c))\int_{x-ct}^{x+ct}g(s)ds\)。
(2) 当 \(g=0\)，积分项消失，解为 \(y=f(x-ct)/2+f(x+ct)/2\)。其中第一项以波速 \(c\) 向正 \(x\) 方向传播，第二项以同速向负方向传播，各具有初始位移峰值的一半；两支叠加在 \(t=0\) 正好复现原位移。
(3) 若 \(f=0\)，则 \(y(x,t)=(1/(2c))\int_{x-ct}^{x+ct}g(s)ds\)。对支持集位于 \([a,b]\) 的 \(g\)，仅当区间 \([x-ct,x+ct]\) 与 \([a,b]\) 相交时扰动可非零，即只能传播到初始支持集外扩距离 \(ct\) 的区域。
(4) 动能与应变能总和 \(E=\frac12\int_{-\infty}^{\infty}[\mu y_t^2+T y_x^2]dx\)，其中 \(c^2=T/\mu\)。对时间求导并利用波方程，\(\dot E=[Ty_xy_t]_{-\infty}^{\infty}=0\)，只要远处扰动充分衰减。这个积分恒等式与有限特征传播速度相容。
